% ---------------------------------------------------------------------------
% 9 Related work
% Sources verified in reviews/2026-08-08-paper-sources.md.
% ---------------------------------------------------------------------------

\section{Related work}

\subsection{Private transformer inference}

Private neural-network inference has long combined cryptographic tools according to the
operation being evaluated. Gazelle established an influential hybrid design in which
homomorphic encryption handles linear layers and garbled circuits handle
nonlinearities~\cite{gazelle2018}. Transformer systems inherited this decomposition but had
to contend with attention, normalization, and activation functions at substantially greater
scale. Iron, Nimbus, and BumbleBee combine homomorphic encryption with secret sharing,
oblivious transfer, and specialized two-party protocols for transformer
inference~\cite{iron2022,nimbus2024,bumblebee2025}. Their protocols optimize a joint secure
computation between parties; the present work instead keeps the server's linear path in CKKS
and assigns exact nonlinear evaluation to the data owner at explicit boundaries. This choice
changes both the security contract and the allocation of work between the parties.

Safhire is the closest architectural antecedent~\cite{safhire2025}. It likewise places
homomorphically evaluated linear layers on the server, then has the client decrypt an
intermediate, apply the nonlinearity in plaintext, and return a fresh encryption. Safhire
studies convolutional networks and adds randomized output permutation to make reconstruction
from released intermediates harder. We therefore do not claim the client-assisted pattern as
new. The contribution here is to test that pattern on a released genomic transformer at its
task-defined prompt, retain the released weights and exact nonlinearities, verify the
encrypted block through an independently checked plaintext reference, and expose the resulting
arithmetic, boundary schedule, and memory behavior. This evaluation applies the established
protocol pattern to a released genomic transformer with a reproducible verification chain.

THOR demonstrates full non-interactive transformer evaluation under homomorphic encryption using
a circuit and packing strategy designed for that objective, including diagonal-major organization
and compact packing~\cite{thor}. Its result establishes that the memory wall observed here belongs
to the evaluated backend, parameters, packing, and circuit. Client assistance addresses a distinct
engineering question: whether released computation can be preserved when the data owner is
available to evaluate exact nonlinearities at declared boundaries.

\subsection{Homomorphic encryption for genomics}

Homomorphic encryption already has a substantial genomic-computing lineage. The iDASH secure
genome analysis competition evaluated encrypted association studies, sequence search, and
related privacy-preserving workflows~\cite{idash2018}. Private genomic-query systems combine
homomorphic encryption with hashing and set-intersection techniques to test for variants
without disclosing either the stored genome or the query~\cite{privategenomicqueries2017}.
Those workloads center on searches or prescribed statistics over genomic records. DNAGPT inference
instead composes dense projections, attention, normalization, and a feed-forward network while
preserving the numerical behavior of a released foundation model. This study extends the genomic
HE lineage to the feasibility boundary of that transformer forward pass.

\subsection{GPU-accelerated homomorphic encryption}

GPU CKKS libraries move number-theoretic transforms, key switching, rotations, and
ciphertext arithmetic onto accelerators; the evaluated backend supplies those primitives and
interoperates with the host cryptographic runtime~\cite{gpu-ckks-backend,openfhe}. Its API
also shaped the available systems designs. In particular, the absence of ciphertext
serialization prevented process-sharded evaluation, leaving thread-level concurrency and
in-process device placement as the implementable paths for this implementation.

Recent multi-GPU transformer work treats placement and communication--computation overlap as
first-class scheduling problems~\cite{aegis}. That literature motivates the placement and overlap
options discussed in Section~\ref{sec:optimization}. The current implementation neither serializes
ciphertexts between processes nor exercises a networked client/server transport, so those paths
must be evaluated directly before adopting a multi-device schedule.

\subsection{Homomorphic matrix primitives}

Matrix layout determines which rotations, plaintext encodings, and reductions an encrypted
linear layer requires. THOR's diagonal-major representation and compact packing specialize
that layout for transformer inference~\cite{thor}; HE-BLAS develops a broader library of
homomorphic linear-algebra kernels and shape-aware algorithms~\cite{heblas2025}. These works
motivate the remaining targets in Section~\ref{sec:optimization}: reducing repeated packing
work, copying fewer distinct diagonal plaintexts, and replacing segmented slot reductions
with kernels matched to the actual tensor shapes. Packing, key material, level consumption, and
the operation schedule are coupled. Integrating a candidate primitive therefore includes
revalidation against the frozen plaintext reference and checked schedule before performance is
compared.
